\scanpage{103}
% Archive pencil foliation: 102.

Lorsque $k$ est de caractéristique nulle, on a toujours un foncteur fibre naturel: le \uline{foncteur de Hodge,} qui à la cohomologie motivique d'une variété (projective lisse) $X$ associe le vectoriel bigradué $\handwritten{\coprod}H^q(X,\handwritten{\underline{\Omega}}^p_{X/k})$ Donc pour le groupe de Galois motivique correspondant $G$, on trouve un homomorphisme naturel\quad\handwritten{\tikz[baseline=0pt,x=1em,y=1em]{\draw (0,0)--(.12,.12);}}
\diagram{\begin{tikzcd}[cramped]\underline{G}_m^2\arrow[r]&G,\end{tikzcd}}
i.e. deux homomorphismes commutant l'un à l'autre\hspace{1em}\handwritten{\tikz[baseline=0pt,x=1em,y=1em]{\draw (.06,-.15)--(.1,-.32);\fill (0,-.4) circle (.045);}}
\diagram{$i_1,i_2:$\begin{tikzcd}[cramped]\underline{G}_m\arrow[r]&G.\end{tikzcd}}
Le fait que la graduation totale dans la cohomologie de Hodge corresponde au poids des motifs s'exprime par la relation
\[
i(\handwritten{\lambda})=i_1(\handwritten{\lambda})\,i_2(\handwritten{\lambda});
\]
on fera attention que $i_1$ et $i_2$ ne sont pas centraux (car la bigraduation en cohomologie de Hodge ne correspond pas à une bigraduation d'un motif !)

2) Prenant toujours pour $\underline{C}$ la c\unsure{\deleted{t}a|\deleted{a}t}\unsure{\deleted{a}t|\deleted{t}a}\unsure{\deleted{g}é|\deleted{é}g}gorie des motifs sur $k$, \unsure{a|e}\unsure{\deleted{m}v|\deleted{n}v|\deleted{v}m|\deleted{v}n}ec $k$ de car. nulle, on a un autre foncteur fibre canonique, le \uline{foncteur de De Rha\unsure{m|n}}\rlap{\handwritten{\tikz[baseline=0pt,x=1em,y=1em]{\draw (0,-.7)--(.3,-.55);}}} qui associe à la cohomologie motivique d'une variété $X$ le vectoriel $H^*(X,\handwritten{\Omega}^*_{X/k})$ (espace d'hypercohomologie). Ce vectoriel n'est plus bigradué\rlap{\handwritten{\tikz[baseline=0pt,x=1em,y=1em]{\draw (0,-.65)--(.2,-.5);}}}\handwritten{\unsure{é,|en,|ée,}} mais seulement gradué et filtré,\unsure{\deleted{x}\handwritten{(}|\deleted{\handwritten{(}}x}la filtration étant celle associée à la suite spectrale d'hypercohomologie, commençant avec la cohomologie de Hodge\unsure{\deleted{x}\handwritten{)}|\deleted{t}\handwritten{)}|\deleted{\handwritten{)}}x|\deleted{\handwritten{)}}t}. On sait d'ailleurs que cette suite spectrale dégénère (théorie de Hodge), donc $\operatorname{Gr}(H_{\mathrm{DR}}(V))\handwritten{\simeq}H_{\mathrm{Hdg}}(V)$ (isomorphisme fonctoriel en le motif $V$\inserted{\rlap{\handwritten{\begin{tikzpicture}[baseline=-.3ex,x=1cm,y=1cm]
\draw (.25,.7) .. controls (-.45,.5) and (-.5,.02) .. (0,-.05) -- (1.6,-.05) .. controls (1.9,.05) and (1.95,-.1) .. (2,-.2) .. controls (2.1,-.05) and (2.2,-.05) .. (2.4,-.05) -- (6.8,-.2) .. controls (7.6,-.25) and (7.85,.1) .. (7.65,.5) .. controls (7.5,.7) and (7.2,.7) .. (7.05,.65);
\draw (2,-.2) .. controls (2.1,-.4) and (1.95,-.65) .. (2,-.9);
\end{tikzpicture}}}\handwritten{où on gradue le $2^{\unsure{\mathrm{em}|\mathrm{eme}}}$ \unsure{membre|membr} p\uline{ar $p$}}}). On peut se proposer d'analyser à quelle structure supplémentaire, sur le groupe de Galois motivique associé au foncteur fibre $H_{\mathrm{DR}}$, correspond la filtration canonique de ce foncteur.


De façon générale, étant donné une catégorie tensorielle $C$ sur $k$ \handwritten{\tikz[baseline=0pt,x=1em,y=1em]{\draw (0,.8)--(.15,1.02);}}munie d'un foncteur fibre $F$, on peut se proposer de déterminer les filtra-

